\subsection{Repository Structure}\label{repository-structure}

GitHub was used as the team's code repository management platform. From the root of
\href{https://github.com/Eines-Informatiques-Avancades/Project_2026_II}{the project},
the code is organized into three main folders: \texttt{src}, \texttt{results}, and \texttt{docs}.

\begin{itemize}
  \item \texttt{src} contains:
  \begin{itemize}
    \item the Makefile and shell scripts,
    \item the different versions of the main program (serial and parallel),
    \item a \texttt{confs} folder where initial configurations are stored and runtime options are controlled through \texttt{input.dat},
    \item a \texttt{lib} folder containing the Fortran modules, Python analysis scripts, and plotting style dependencies.
  \end{itemize}
  \item \texttt{results} contains the simulation outputs.
  \item \texttt{docs} contains the documentation and the report sources, including the \texttt{img} folder used for embedded figures.
\end{itemize}

In addition to these folders, the repository includes a \texttt{README} file with execution instructions and dependency information. The main simulation code is written in Fortran, while the post-processing and plotting workflow is implemented in Python.

\paragraph{Module library overview.}
The main program files rely on the following supporting modules:

\begin{itemize}
  \item \texttt{energy.f90} and \texttt{energy\_all\_atoms.f90}: implement the two energy models, allowing the code to switch between TraPPE-UA and OPLS-AA-style effective backbone potentials depending on the \texttt{explicit\_h} option read from \texttt{input.dat}.
  \item \texttt{initial-conf.f90}: generates the initial polymer configuration.
  \item \texttt{io.f90}: handles file input and output operations.
  \item \texttt{monte\_carlo.f90}: implements the rotation algorithm and the Monte Carlo update step.
  \item \texttt{observables.f90}: computes the structural observables, including end-to-end distance \(R_{ee}\), radius of gyration \(R_g\), and torsional angles.
  \item \texttt{parameters.f90}: stores global simulation parameters and interfaces with the settings defined in \texttt{input.dat}.
  \item \texttt{plot\_*.py}: post-processing and plotting scripts for observables and serial-versus-parallel comparisons.
  \item \texttt{requirements.txt}: lists the Python dependencies.
  \item \texttt{science.mplstyle}: defines the plotting style used in the generated figures.
\end{itemize}

\subsection{Code Management via GitHub}\label{code-management-via-github}

The project leader reviewed pull requests, checked that incoming changes followed the project guidelines, and merged them into the main branch. Team coordination was handled through a shared WhatsApp group, where members reported progress and informed the rest of the group whenever new code had been merged.

In most cases, integration was straightforward. However, overlapping uploads from different team members occasionally produced merge conflicts. This led the team to adopt a consistent workflow based on stashing and rebasing before recommitting local work. The process used was the following:

\begin{enumerate}
  \item Stash local uncommitted changes:
\begin{verbatim}
git stash
\end{verbatim}

  \item Pull the latest version of the main branch:
\begin{verbatim}
git checkout master
git pull main master
\end{verbatim}

  \item Rebase the local feature branch onto the updated main branch:
\begin{verbatim}
git checkout <fork-feature>/<fork-branch>
git rebase master
\end{verbatim}

  \item Resolve conflicts and re-apply the stashed changes:
\begin{verbatim}
git stash pop
\end{verbatim}

  \item Commit and push the rebased branch:
\begin{verbatim}
git push <fork-remote> <fork-branch> --force
\end{verbatim}
\end{enumerate}

This workflow helped keep the repository history cleaner and reduced repeated integration issues during collaborative development.

\subsection{Makefile}\label{makefile}

As the project evolved, the Makefile also became more structured. It was used to automate compilation, execution, and post-processing for both the serial and parallel versions of the code.

The main compilation variables are:

\begin{verbatim}
FC        = gfortran
FFLAGS    = -O2 -Wall -Wextra
OBJ_DIR   = ../bin
EXE       = $(OBJ_DIR)/main_serial.x
LIB_OBJ   = $(OBJ_DIR)/parameters.o \
            $(OBJ_DIR)/io.o \
            ...
MAIN_OBJ  = $(OBJ_DIR)/main_serial.o

NP          ?= 4
OMP_THREADS ?= 2
\end{verbatim}

Their roles are summarized below:

\begin{itemize}
  \item \texttt{FC}: the Fortran compiler, initially set to \texttt{gfortran}.
  \item \texttt{FFLAGS}: compiler flags; \texttt{-O2} enables optimization, while \texttt{-Wall -Wextra} activates additional warnings.
  \item \texttt{OBJ\_DIR} and \texttt{EXE}: define the directory for compiled objects and the name of the executable.
  \item \texttt{LIB\_OBJ} and \texttt{MAIN\_OBJ}: list the object files needed for the library modules and the main program.
  \item \texttt{NP} and \texttt{OMP\_THREADS}: define the number of MPI processes and OpenMP threads used in parallel executions. These values can be overridden from the command line, for example with \texttt{make run\_parallel NP=7 OMP\_THREADS=3}.
\end{itemize}

To support both serial and parallel builds, the Makefile checks whether the requested target contains the word \texttt{parallel}. If so, it verifies that the MPI compiler wrapper is available and then switches the compiler from \texttt{gfortran} to \texttt{mpif90}.

\begin{verbatim}
ifneq (, $(findstring parallel,$(MAKECMDGOALS)))
  MPI_BIN := $(shell command -v mpif90 2>/dev/null)
  ifeq ($(strip $(MPI_BIN)),)
    $(error "MPI Compiler is required but 'mpif90' was not found!")
  endif
  FC = mpif90
  ...
endif
\end{verbatim}

Instead of writing an explicit build rule for every \texttt{.f90} file, a pattern rule is used:

\begin{verbatim}
$(OBJ_DIR)/%.o: lib/%.f90
    @mkdir -p $(OBJ_DIR)
    $(FC) $(FFLAGS) -J$(OBJ_DIR) -I$(OBJ_DIR) -c $< -o $@
\end{verbatim}

This rule states that any object file in \texttt{bin/} can be built from its corresponding source file in \texttt{lib/}. The \texttt{-c} option compiles without linking, while \texttt{-J} and \texttt{-I} direct the compiler to store and search module files in the object directory.

The final executable is linked through the default target:

\begin{verbatim}
all: $(EXE)

$(EXE): $(LIB_OBJ) $(MAIN_OBJ)
    @mkdir -p $(OBJ_DIR)
    $(FC) $(FFLAGS) -o $@ $(LIB_OBJ) $(MAIN_OBJ)
\end{verbatim}

Finally, convenience targets are declared as \texttt{.PHONY} rules:

\begin{verbatim}
.PHONY: all clean run figures pipeline

run: all
    @mkdir -p ../results
    $(EXE)

figures:
    python lib/plot_results.py

clean:
    rm -rf $(OBJ_DIR)/*.o $(OBJ_DIR)/*.mod $(EXE)

pipeline:
    $(MAKE) run
    $(MAKE) figures
\end{verbatim}

These shortcuts make it possible to compile, execute, clean, and post-process the project with short commands such as \texttt{make run}, \texttt{make clean}, and \texttt{make pipeline}.

\subsection{Shell Scripting}\label{shell-scripting}
% Itxaso bit added + Arthur's explanation merged

Shell scripting was used to make our runs compatible with the CERQT2
cluster environment and to automate benchmarking. The scripts
initialize the runtime environment, load the required modules, and set
variables such as \texttt{OMP\_NUM\_THREADS}, \texttt{NSLOTS}, and
\texttt{MPLBACKEND=Agg} for stable non-interactive execution. They were
also used to automate parameter sweeps, build explicit trajectory
manifests for the observables analysis, and collect timing data into
summary files for MPI and OpenMP tests.

A qsub script \texttt{run\_collab.qsub} was also written during the
testing of the program to submit an individual portion of the simulation
to the HPC cluster. It involved loading the appropriate modules for that
queue, selecting a custom makefile (separate from the main program's
Makefile), indentifying the correct commandline arguments to pass at
compile time, and submit the mpirun command.